% Capítulo 5 (Metodologia) — Slide — Regimes de demanda na taxonomia de Syntetos-Boylan
\begin{frame}{Regimes de demanda na taxonomia de Syntetos-Boylan}
  \centering
  \resizebox{0.6\linewidth}{!}{%
  \begin{tikzpicture}[
      every node/.style={font=\tiny\sffamily},
      regime/.style={rectangle, rounded corners=5pt, draw=slategray!55, fill=cardbg,
                     minimum width=4.3cm, minimum height=1.95cm, align=center, text width=3.95cm},
      axis/.style={->, thick, slategray}
    ]
    \node[regime] (erratic) {%
      \textbf{\scriptsize Erratic}\\[1pt]
      Baixa intermitência + alta variabilidade\\
      Frequente, tamanhos irregulares. Desafio: variabilidade.\\[2pt]
      \spark{0/0.3,1/0.9,2/0.5,3/0.95,4/0.25,5/0.8,6/0.4,7/0.85}};
    \node[regime, right=0.9cm of erratic] (lumpy) {%
      \textbf{\scriptsize Lumpy}\\[1pt]
      Alta intermitência + alta variabilidade\\
      Muitos ciclos sem demanda e picos irregulares. O mais desafiador.\\[2pt]
      \spark{0/0,1/0,2/0.9,3/0,4/0.3,5/0,6/0,7/0.95}};
    \node[regime, below=0.6cm of erratic] (smooth) {%
      \textbf{\scriptsize Smooth}\\[1pt]
      Baixa intermitência + baixa variabilidade\\
      Frequente e estável. Mais fácil de prever e repor.\\[2pt]
      \spark{0/0.7,1/0.65,2/0.75,3/0.7,4/0.68,5/0.72,6/0.7,7/0.69}};
    \node[regime, below=0.6cm of lumpy] (intermittent) {%
      \textbf{\scriptsize Intermittent}\\[1pt]
      Alta intermitência + baixa variabilidade\\
      Muitos ciclos sem demanda; tamanho relativamente estável quando ocorre.\\[2pt]
      \spark{0/0,1/0,2/0.6,3/0,4/0,5/0.6,6/0,7/0.6}};

    \draw[axis] ($(smooth.south west)+(-0.15,-0.4)$) -- ($(intermittent.south east)+(0.15,-0.4)$)
      node[pos=0, below]{\tiny baixa}
      node[pos=1, below]{\tiny alta}
      node[midway, below=9pt]{\tiny Intermitência (ADI)};
    \draw[axis] ($(smooth.south west)+(-0.4,-0.15)$) -- ($(erratic.north west)+(-0.4,0.15)$)
      node[pos=0, left]{\tiny baixa}
      node[pos=1, left]{\tiny alta}
      node[midway, left=9pt, rotate=90]{\tiny Variabilidade (CV$^2$)};
  \end{tikzpicture}}
  \vspace{0.25em}

  \small\keymessage{Lumpy combina muitos ciclos sem demanda com alta variabilidade nos tamanhos dos pedidos, tornando políticas baseadas em médias globais menos robustas.}
  \note{
    Antes de detalhar os Perfis Operacionais de Demanda, vale fixar a
    leitura do plano ADI por CV$^2$ que fundamenta a taxonomia clássica
    usada como ponto de partida da caracterização. O eixo horizontal mede
    intermitência pelo ADI -- intervalo médio entre demandas positivas; o
    eixo vertical mede variabilidade pelo CV$^2$ -- a variância normalizada
    do tamanho da demanda. Combinando os dois, a taxonomia de
    Syntetos-Boylan define quatro regimes. Smooth é demanda frequente e
    estável, mais fácil de prever e repor. Erratic é frequente, mas com
    tamanhos irregulares -- aqui o desafio é a variabilidade, não a
    ocorrência. Intermittent tem muitos ciclos sem venda, mas quando a
    demanda ocorre o tamanho é relativamente estável -- o desafio é a
    ocorrência esporádica. Lumpy combina os dois problemas: muitos ciclos
    sem demanda e, quando ocorre, picos irregulares de magnitude -- é o
    regime mais desafiador para previsão e reposição. Primeiro eu defino
    como a literatura classifica regimes de demanda; o POD, no próximo
    slide, vai além dessas duas variáveis; e mais adiante, no Capítulo de
    Resultados, mostro onde o recorte avaliado cai nesse mesmo plano.
  }
\end{frame}
